% Français 9115, batch 14 — Gallica views 261–280.
% Pass: Opus 5 (claude-opus-5), 2026-09-19 — first pass, unchecked against
% the views by a human.
%
% What the twenty views hold. Not the Euler copy. The Latin fair copy that
% filled batches 8–10 has stopped before this batch: what these leaves carry
% is a continuous piece of Sophie Germain's own French prose and algebra on
% the vibrating elastic lamina held by a stylet at an interior point — her
% own memoir, argued against Euler's « Investigatio motuum » (which she
% cites by its section numbers), against Chladni's « traité d'acoustique »
% and against a memoir of « Mr Ricati ». It is a worked draft, not a fair
% copy: struck words, words written over words, additions between the lines
% and four notes of her own in the margins.
%
% Only ten of the twenty views carry writing, and they are the even ones.
% Each leaf of this run is written on one side only and photographed twice:
% views 262, 264 … 280 are the written sides, and views 261, 263 … 279 are
% the blank other sides, on which the writing of the leaf shows through
% reversed. This was checked by mirroring the crop (\texttt{magick … -flop})
% and laying it on the written view: view 279 flopped is view 278 (page 136)
% line for line, and view 261 flopped is the page numbered 127, which is
% view 260 and belongs to batch 13. The ten blank views are therefore not
% transcribed and no placeholder stands for them. There are no repeats and
% no microfilm targets in this batch.
%
% The first line of view 262 continues a sentence begun on view 260, which
% is not in this batch; the show-through of view 261 identifies that page as
% the one numbered 127, so the seam is known even though the page itself was
% not read here.
%
%   v. 262  (p. 128) the end of the stylet calculation: the transcendental
%           equation in ω, then λ = 1/2 as its solution, which makes the
%           lamina move as two half-length laminae each with one end on the
%           stylet; and that the two tangent equations cannot both agree
%           with the end conditions in any other case.
%   v. 264  (129) § 13: why the lamina may be treated as two separate
%           laminae when the stylet is at its middle and in no other case.
%           Two notes of the writer's in the left margin, keyed « + » and
%           « * »; the first is cancelled by a long diagonal stroke.
%   v. 266  (130) the conclusion: for laminae giving the same sound the part
%           between the appuiée end and the first node is the same length,
%           whatever the length of the lamina; what remains to be proved.
%   v. 268  (131) § 14: the proof. Euler's general y for a point L at λa,
%           the condition y = 0 there, the ordinates at λ ∓ λ′, their sum
%           zero, and the condition it reduces to.
%   v. 270  (132) the same with the coefficients of the other case; λ = 1/2
%           is the only value that satisfies it; a node can never stand at
%           equal distances from the two nodes next to it. § 15 opens the
%           case of a lamina with both ends simply appuiées.
%   v. 272  (133) the parts between successive nodes are unequal: Chladni
%           « traité d'acoustique pag 4 §§ 71 » is said to have made them
%           double; Riccati, following Euler, computed values which give the
%           inequality but drew no conclusion from it. Two lines in a much
%           finer pen at the foot, with the reference « V. §§ XXX p. 503 »,
%           the paragraph number left as three crosses.
%   v. 274  (134) the three sine conditions for a node at the stylet's
%           point, and § 16 on the analogy between the equations of the
%           several cases. A note of the writer's in the left margin.
%   v. 276  (135) the six cases of Euler's memoir set out in a list with
%           their section numbers — 9, 23, 30, 35, 39, 43 — and the equation
%           of ω for each; then the equations that hold when the stylet is
%           applied, for both ends appuiées, both libres, both fixées. The
%           last of these is the formula that opens view 262.
%   v. 278  (136) § 17: the case of one end fixed and the other free,
%           referred to « § 47 et suiv. du memoire d'Euler »; each equation
%           read as the product of two conditions, one for each portion of
%           the lamina.
%   v. 280  (137) § 18: despite the analogies the method yields nothing
%           applicable, there being no general solution except when both
%           ends are appuiées. A long note of the writer's in the left
%           margin says the same in other words.
%
% The numbers on the leaves, and why no \folio{} is written. Two series,
% both in ink, one per written leaf; the blank views carry none.
% (1) Top right: 128 (v. 262), 129 (264), 130 (266), 131 (268), 132 (270),
% 133 (272), 134 (274), 135 (276), 136 (278), 137 (280) — without a gap,
% continuing the ink series recorded in batches 8 to 10.
% (2) Top centre, underlined, several of them written over a first figure:
% 21 (v. 262), 22 (264), 23 (266), 24 (268), 25 (270), 25 (272), 26 (274),
% 27 (276), 28 (278), 29 (280). The two figures read 25 stand one after the
% other, which suggests one of them was corrected; their second digit is
% repassed and the reading is not certain. This is the writer's own
% pagination of the piece, not a foliation.
% Besides these, the paragraph numbers of the text stand alone in the left
% margin: 13 (v. 264), 14 (268), 15 (270), 16 (274), 17 (278), 18 (280) —
% the last three in large blotted figures. View 272 carries a further small
% mark at the top left, read « 2.eme » with doubt.
% No pencilled foliation of the Bibliothèque was read on any view of this
% batch, so no \folio{} is written, as in batches 1 to 10. Every number
% above was read on its view; none is inferred.
%
% Notation. Hers and Euler's together: ω for the angle a/f of the memoir,
% λ for the fraction EL/EF and λ′ for a distance measured from L, a the
% length of the lamina, and the exponentials e^{λω}, e^{(1−λ)ω},
% e^{(λ−1)ω} as she writes them. The double signs are written in two tiers,
% a bar under or over the cross, and are set \pm and \mp as they read; the
% superior sign belongs to one state of the ends, the inferior to the other,
% as the leaves say. « tang. », « sin. » and « cos. » are written with and
% without their stop in the same line and the pass has not tried to record
% every stop. Lines of dots opening a formula are hers and stand for the
% factor she leaves out, as at views 186 and 189 of batch 10; they are set
% \ldots. Her spelling is kept: « egals » and « égals » beside « egales »,
% « fesant », « d'avantage », « Apresent », « parceque », « adire »,
% « precedans », « differens », « Ricati ». Words added between the lines
% are put at the place the leaf gives them and the addition is said in a
% note; words struck are kept in \struck{}. The « + », « * » and « × » in
% the text are her keys to the marginal notes, and are set as signs.
%
% What could not be read. \ill{} stands in five places: a word heavily
% cancelled after « la methode d'Euler » (v. 274) and another among the
% interlinear words above « determinent » on the same view; a sign under a
% thick rature in the first formula of view 268; a cancelled word after
% « dans le cas ou » (v. 278); the end of « montrer q… » where the paper is
% wanting (v. 280). Readings offered with doubt are marked \uncertain{}:
% the sign read « et » in the interlinear addition of view 264 and the two
% « \&c. » of that view; « car cette position est » in its first marginal
% note; « cos » under a blot at view 268; « V. » before « traité
% d'acoustique » and the page « 503 » at view 272; « problemes » in the
% marginal note of view 274; « dedu » and the section number « 9 » at view
% 280. Nothing was guessed to fill a gap.
%
% Nothing here was checked against any published transcription or edition of
% Germain, of Euler, of Chladni or of Riccati.
\documentclass[11pt,a4paper]{article}
% The preamble shared by every transcription of the mathematicians' manuscripts in Gallica.
%
% It rests on one idea: **uncertainty compounds, it does not resolve.** A
% transcription that smooths over a doubtful word has destroyed the only thing
% it was made for — a reader must be able to tell, at every point, what was on
% the page and what was supplied. The apparatus macros below are therefore the
% critical apparatus, not decoration.
%
% The same commands are recognised by `scripts/render.mjs`, which produces the
% site's reading view, and by `scripts/tei.mjs`. A macro added here must be
% added there too, or rendering fails loudly — which is the wanted behaviour.
%
% Comments here are English, like the rest of the repository. The documents
% this preamble sets are French, and that is deliberate: see the skills.

\ProvidesPackage{germain}[2026/09/15 Transcriptions of mathematicians' manuscripts in Gallica]

\usepackage{amsmath,amssymb,amsthm}
\usepackage{mathrsfs}
% Kept from the parent preamble so the renderer's diagram support stays
% available; these manuscripts draw no commutative diagrams, and nothing here uses it.
\usepackage{tikz-cd}
\usepackage[english,french]{babel}
\usepackage{fontspec}

% --- Greek ------------------------------------------------------------------
% The text face stays fontspec's default, Latin Modern, which has no Greek:
% XeTeX drops every glyph it lacks with a line in the log and nothing on the
% page, and the Greek words the manuscripts quote vanished from the PDFs. So
% the Greek and Greek Extended blocks get a XeTeX character class of their own,
% and entering or leaving a run of it switches the family to CMU Serif — the
% same Computer Modern design, with polytonic Greek — and back. Only the
% family changes: a Greek word in \textit or \textbf keeps its shape and series.
% The transitions to and from every other class carry the tokens that class
% has towards an ordinary letter, so French spacing before « ; » or inside
% « » still holds after a Greek word. No grouping, so no brace or \endgroup
% can come out unbalanced; maths is left alone, since XeTeX inserts nothing
% there. Kept identical in `exercices.sty`.
\makeatletter
\newfontfamily\germainGreekfont{cmun}[
  Extension=.otf, NFSSFamily=germaingreek,
  UprightFont=*rm, ItalicFont=*ti, SlantedFont=*sl,
  BoldFont=*bx, BoldItalicFont=*bi, BoldSlantedFont=*bl]
\newXeTeXintercharclass\germain@greek
\def\germain@greekrange#1#2{%
  \@tempcnta=#1\relax
  \loop
    \XeTeXcharclass\@tempcnta=\germain@greek
    \ifnum\@tempcnta<#2\advance\@tempcnta\@ne\repeat}
\germain@greekrange{"0370}{"03FF}
\germain@greekrange{"1F00}{"1FFF}
\def\germain@greekprev{\rmdefault}
\def\germain@greekin{%
  \edef\germain@greekprev{\f@family}\fontfamily{germaingreek}\selectfont}
\def\germain@greekout{\fontfamily{\germain@greekprev}\selectfont}
\def\germain@greekjoin{%
  \XeTeXinterchartokenstate=\@ne
  \@tempcnta=\z@
  \loop
    \ifnum\@tempcnta=\germain@greek\else
      \XeTeXinterchartoks\germain@greek\@tempcnta=\expandafter{%
        \expandafter\germain@greekout\the\XeTeXinterchartoks\z@\@tempcnta}%
      \XeTeXinterchartoks\@tempcnta\germain@greek=\expandafter{%
        \the\XeTeXinterchartoks\@tempcnta\z@\germain@greekin}%
    \fi
    \ifnum\@tempcnta<4095\advance\@tempcnta\@ne\repeat}
% babel-french sets its own classes, and saves and restores the interchar
% state, each time a language is selected: join again after it does.
\addto\extrasfrench{\germain@greekjoin}
\addto\extrasenglish{\germain@greekjoin}
\AtBeginDocument{\germain@greekjoin}
\makeatother

\usepackage{geometry}
\geometry{a4paper,margin=25mm,marginparwidth=28mm}
\usepackage{marginnote}
\usepackage{xcolor}
\usepackage[normalem]{ulem}
\setlength{\emergencystretch}{3em}
\usepackage{hyperref}
\hypersetup{colorlinks=true,linkcolor=black,urlcolor=black,pdfborder={0 0 0}}

% --- Batch metadata ---------------------------------------------------------
% Taken as they stand by the rendering script, which shows them at the head of
% the reading view. They fix what the file claims to cover. The macro names are
% the parent project's, so that the scripts stay shared; \folder here means the
% volume's slug — `fr-9115` — and \pages counts Gallica views.

\newcommand{\folder}[1]{\gdef\grothFolder{#1}}
\newcommand{\batch}[1]{\gdef\grothBatch{#1}}
\newcommand{\pages}[2]{\gdef\grothFirst{#1}\gdef\grothLast{#2}}
\newcommand{\foldertitle}[1]{\gdef\grothFoldertitle{#1}}

% The holder's own shelfmark, printed as they write it: « Français 9115 ».
\newcommand{\shelfmark}[1]{\gdef\grothShelfmark{#1}}

% The dating, copied verbatim from the holder's catalogue — never inferred
% here. « XIXe siècle » is the BnF's; a pass that dates a leaf from its content
% says so in a \note{}, as its own inference.
\newcommand{\dating}[1]{\gdef\grothDating{#1}}

% The status notice, printed in the title block and repeated as a diagonal
% watermark on every page of the PDF. The manuscripts are in the public
% domain; what this line declares is not a right but a quality — an
% unverified first pass by a machine. The skills write the line; a file
% without it gets no watermark, so leaving it out is visible in the output.
\newcommand{\watermark}[1]{\gdef\grothWatermark{#1}}

\gdef\grothFolder{?}\gdef\grothBatch{}\gdef\grothFirst{?}\gdef\grothLast{?}%
\gdef\grothFoldertitle{}\gdef\grothShelfmark{}\gdef\grothDating{}\gdef\grothWatermark{}

% --- The résumé -------------------------------------------------------------
\newenvironment{resume}
  {\small\setlength{\parskip}{0.45em}}
  {\par\medskip\hrule\medskip}

% --- Keywords ---------------------------------------------------------------
% In English, closing the modernised reading's résumé: the modern vocabulary
% under which to search for what the leaves construct. The manifest extracts
% them to tag the volume — this line is the single source of the tags.
\newcommand{\keywords}[1]{\par\smallskip\noindent
  {\footnotesize\textsc{Keywords} --- \textit{#1}}\par}

% --- The critical apparatus -------------------------------------------------

% The Gallica view number, in the margin. This is the anchor that lets a
% disagreement over a formula be settled by looking at *one* image rather than
% a volume of 750. The site uses it to turn the facsimile as the transcription
% is scrolled. A view is one scanned image: usually one side of a leaf.
\newcommand{\page}[1]{%
  \marginnote{\footnotesize\color{gray!70}\textsf{v.\,#1}}%
  \label{page:#1}\ignorespaces}

% The BnF's pencilled foliation, where the leaf carries it and a pass can read
% it: « 198r ». This is what the literature cites — « ff. 198r–208v » — and
% the one line that turns a citation into a view. Recorded, never inferred:
% a folio a pass cannot read is not written.
\newcommand{\folio}[1]{%
  \marginnote{\footnotesize\color{gray!50}\textsf{f.\,#1}}\ignorespaces}

% The modernised reading's coarser counterpart to \page{}: which views a
% section is drawn from.
\newcommand{\pagerange}[2]{%
  \marginnote{\footnotesize\color{gray!70}\textsf{v.\,#1--#2}}%
  \ignorespaces}

% Illegible: never guessed.
\newcommand{\ill}{\textcolor{red!60!black}{[\,\ldots\,]}}

% A doubtful reading: the word is given, and flagged as uncertain.
\newcommand{\uncertain}[1]{\uline{#1}}

% An editorial addition — a missing letter, an implied word.
\newcommand{\add}[1]{\textcolor{blue!50!black}{[#1]}}

% Struck out by the writer. What was crossed out is part of the document.
\newcommand{\struck}[1]{\sout{#1}}

% The transcriber's note, on the state of the page or the mathematics.
\newcommand{\note}[1]{\par\smallskip{\small\color{gray!60!black}\textit{[#1]}}\par\smallskip}

% A marginal note of hers, distinct from the body of the text.
\newcommand{\marginal}[1]{\par\smallskip\noindent\hspace{1em}%
  {\small\color{gray!60!black}#1}\par\smallskip}

% --- Title ------------------------------------------------------------------
% The title block is French, like the documents themselves.

\AddToHook{shipout/background}{%
  \ifx\grothWatermark\empty\else
    \begin{tikzpicture}[remember picture, overlay]
      \node[rotate=52, scale=3.4, align=center, text=gray!18,
            font=\sffamily\bfseries]
        at (current page.center) {\grothWatermark};
    \end{tikzpicture}%
  \fi}

\AtBeginDocument{%
  \begin{center}
    {\large\bfseries Manuscrits de mathématiciens (Gallica) ---
      \ifx\grothShelfmark\empty\grothFolder\else\grothShelfmark\fi}\\[2pt]
    {\itshape\grothFoldertitle}\\[4pt]
    {\small vues \grothFirst--\grothLast
      \ifx\grothBatch\empty\else\ (lot \grothBatch)\fi}\\[2pt]
    \ifx\grothDating\empty\else
      {\small\color{gray!60!black}Datation du catalogue : \grothDating}\\[2pt]
    \fi
    \ifx\grothWatermark\empty\else
      {\small\bfseries\color{red!45!gray}\grothWatermark}\\[2pt]
    \fi
    {\footnotesize\color{gray!60!black}Transcription établie d'après les images IIIF de Gallica.
     Source gallica.bnf.fr / Bibliothèque nationale de France.\\
     Les lectures incertaines sont soulignées, les illisibles notées [\,\ldots\,],
     les ajouts entre crochets ; \textsf{v.} numérote les vues de Gallica,
     \textsf{f.} la foliotation au crayon de la BnF.}
  \end{center}
  \vspace{1em}}

\endinput


\folder{fr-9115}
\shelfmark{Français 9115}
\batch{14}
\pages{261}{280}
\dating{1801-1900}
\watermark{Lecture automatique\\non vérifiée}
\foldertitle{Papiers de Sophie GERMAIN. — Recueil de dissertations et problèmes mathématiques et physiques.}

\begin{document}

\note{Numérotation des feuillets. Chaque feuillet de ce lot n'est écrit que
d'un côté : les vues paires portent le texte, les vues impaires sont le
revers blanc du même feuillet, où l'écriture transparaît à l'envers ; elles
ne sont pas transcrites. Deux séries de nombres à l'encre, une par feuillet
écrit. En haut à droite : 128 (vue 262), 129 (264), 130 (266), 131 (268),
132 (270), 133 (272), 134 (274), 135 (276), 136 (278), 137 (280), sans
lacune, à la suite de la série des lots précédents. Au milieu du haut,
souligné et souvent repassé sur un premier chiffre : 21 (262), 22 (264),
23 (266), 24 (268), 25 (270), 25 (272), 26 (274), 27 (276), 28 (278),
29 (280) ; deux vues de suite se lisent 25 et le second chiffre de ces
nombres n'est pas assuré. C'est la pagination de l'auteur, non une
foliotation. Les numéros d'alinéa 13, 14, 15, 16, 17 et 18 sont dans la
marge de gauche des vues 264, 268, 270, 274, 278 et 280. Tous ces nombres
ont été lus sur la vue ; aucun n'est déduit. Aucune foliotation au crayon de
la Bibliothèque n'a été lue sur ces vues, et aucun « f. » n'est donc porté
ici.}

\page{262}

\[ \Big\{\frac{2}{e^{\lambda\omega}+e^{-\lambda\omega}}-\cos.\lambda\omega\Big\}\Big\{\frac{e^{(1-\lambda)\omega}+e^{(\lambda-1)\omega}}{e^{(1-\lambda)\omega}-e^{(\lambda-1)\omega}}\sin(1-\lambda)\omega-\cos(1-\lambda)\omega\Big\} \]
\[ +\Big\{\frac{2}{e^{(1-\lambda)\omega}+e^{(\lambda-1)\omega}}-\cos(1-\lambda)\omega\Big\}\Big\{\frac{e^{\lambda\omega}+e^{-\lambda\omega}}{e^{\lambda\omega}-e^{-\lambda\omega}}\sin.\lambda\omega-\cos.\lambda\omega\Big\}=0 \]
de laquelle il s'agit de tirer les valeurs de l'angle $\omega$

Et on trouve comme dans le cas précédent qu'en fesant $\lambda=\frac{1}{2}$ on a
pour solution de cette equation la condition
\[ \frac{e^{\frac{1}{2}\omega}+e^{-\frac{1}{2}\omega}}{e^{\frac{1}{2}\omega}-e^{-\frac{1}{2}\omega}}\sin.\tfrac{1}{2}\omega-\cos\tfrac{1}{2}\omega=0 \]
qui mène a conclure \struck{que} la lame primitive se meut alors comme si
elle étoit separée en deux autres lames moitiés moins longues qui
auroient chacune une de leurs extremités appuiée au point
d'application du stilet et l'autre fixée.

On trouve également par \struck{le même} calcul déja employé \struck{que} pour le cas
précédent que les equations
\[ \text{tang.}(1-\lambda)\omega=\frac{e^{(1-\lambda)\omega}-e^{(\lambda-1)\omega}}{e^{(1-\lambda)\omega}+e^{(\lambda-1)\omega}} \]
\[ \text{tang.}\lambda\omega=\frac{e^{\lambda\omega}-e^{-\lambda\omega}}{e^{\lambda\omega}+e^{-\lambda\omega}} \]
ne peuvent, dans aucun autre cas que celui de $\lambda=\frac{1}{2}$, s'accorder
l'une et l'autre avec la condition dependante de l'état des extremités.
\note{le feuillet est un morceau de papier plus petit, monté sur une garde : le texte n'occupe que le haut de la vue, et le bas est blanc. Le nombre 128 est à l'encre en haut à droite ; au milieu du haut, un 21 souligné. « primitive » est ajouté au-dessus de la ligne, après « la lame », et mis à sa place. Le mot biffé après « a conclure » est « que », recouvert d'une rature épaisse. Dans la phrase suivante, « le même » est biffé et « déja employé » ajouté au-dessus de « calcul » ; le « que » qui suivait « calcul » est biffé. La première ligne continue la page précédente (vue 260), qui n'est pas de ce lot.}
\page{264}

\marginal{$+$ quelque soit la longueur de la lame dont il s'agit et par
consequent le nombre des nœuds de vibration qui se forment pendant son
mouvement si on suppose qu'elle donne un son \struck{connu} fixe \struck{la}
l'étendue de la partie vibrante comprise entre celle de ses extremités qui
est appuiée d'apres l'\struck{hypothese}, et le nœud de vibration le plus
voisin \struck{est} sera constante \struck{car cette position est}}

\marginal{$*$ si on suppose qu'un son quelconque soit rendu a la fois par
plusieurs lames de differentes longueurs \struck{et} \struck{soit} dont une des
extremités \struck{est} fixe ou libre et l'autre appuiée il se formera dans
chacune d'elles un nombre de nœud dependant de leur longueurs respectives
mais que l'étendue de la partie \struck{comprise} vibrante comprise entre
celles de leurs extremités qui est appuiée et le nœud le plus voisin sera
constante.}

13. Mais, si dans un cas quelconque qui est celui ou le stilet est appuié au
milieu de la lame primitive, la théorie permet de considerer cette lame
comme réellement separée en deux autres lames qui auroient une de leurs
extremités appuiées par l'effet de l'application du stilet, pour comment se
peut-il que la même supposition ne soit plus applicable a aucun autre cas ?

Voici \struck{ce me semble} comment on peut se rendre compte de cette
singularité : \struck{pour que quelque} tout porte a croire que $*$ je suppose
qu'une lame, dont une extremité est appuiée \uncertain{et} l'autre libre ou
fixe, d'une longueur quelconque $*$, rende un son donné en se divisant en
\struck{un certain nombre} plusieurs de parties au moyen d'un nombre impair
de \struck{nœud ou de} \struck{plusieurs} nœuds il semble que l'on peut
concevoir que les \struck{parties intermédiaires soient supprimées pour
laisser subsister que deux separées par un seul nœud, sans que la nature du
son soit changé} deux parties les plus voisines du milieu de la lame soit
supprimées et les portions restantes \struck{continues} contigues, puis
encore les deux \struck{portions} prises de même dans la nouvelle lame ou
quatre dans la première, et ensuite \uncertain{\&c.}
\uncertain{\&c.} jusqu'a ce qu'il n'en reste plus que deux separés par un
seul nœud et que par ces opérations successives le son ne sera pas altéré,
car si on vouloit que \struck{a chaque quelle} le son fut changé, il devroit
l'être a chacune des opérations successives et comme rien ne limite leur
nombre parceque l'on peut concevoir le nombre des nœuds de la lame primitive
infiniment grand, il faudroit supposer que les sons que peuvent rendre
\struck{des portions de lame} les \struck{portions} parties de la lame
situées a ses extremités \struck{sont} suivant l'étendue des lames
auxquelles elles appartiennent infiniment differens entr'eux et
\struck{ce} résultat est contraire a tout ce que l'on sait
\note{le nombre 129 est à l'encre en haut à droite ; au milieu du haut, un 22 souligné, dont le second chiffre est repassé. Le « 13 » qui ouvre la page est dans la marge de gauche, comme un numéro d'alinéa. Les deux notes marginales sont de la main de l'auteur, dans la marge de gauche, et appelées par un « $+$ » et un « $*$ » placés dans le texte ; la première est en outre annulée par un grand trait oblique qui la traverse de haut en bas. « dont une extremité est appuiée $\ldots$ l'autre libre ou fixe » est ajouté au-dessus de la ligne, et mis à sa place ; le signe qui sépare les deux membres est un trait épais, lu « et » avec doute. « plusieurs » est ajouté au-dessus de « un certain nombre » biffé, « nombre impair de » au-dessus de la ligne, « fixe » au-dessus de « connu » dans la note marginale, « sera » au-dessus de « est », « que » au-dessus de la ligne dans la seconde note. Le signe qui ferme « et ensuite » et celui qui ouvre la ligne suivante ont la même forme, une boucle suivie d'un point, lue « \&c. » avec doute. La page se termine sans point ; la phrase se poursuit à la vue 266.}
\page{266}

sur la théorie des corps sonores ; il faut donc conclure que quelque soit
l'étendue de deux lames dont une des extremités est libre ou fixe et l'autre
appuiée, les parties comprises entre l'extremité appuiée, qui est la seule a
considerer ici, et le premier nœud, sont égals entr'elles lorsque les lames
dont il s'agit rendent le même son.

Cela posé, si on considere la lame élastique comme separée en deux parties
par l'application du stilet simplement appuié, il faudra que les deux
\struck{part} portions comprises de part et d'autre entre ce point d'appui
supposé et le premier nœud soient égals entr'elles, car elles
appartiendroient dans cette hypothese a deux lames dont une extremité seroit
libre ou fixe et l'autre appuiée et qui doivent par la nature de la question
rendre le même son. il ne reste donc plus pour eclaircir la difficulté dont
il s'agit, qu'a prouver qu'excepté le cas ou il y a un nœud au milieu de la
longueur de la lame il n'y a aucun point de repos tel que les ordonnées
prises a distance egales de part et d'autre de ce point soient egals
entr'elles et que par conséquent les distances comprises entre le point
d'application du stilet considéré comme simple appui, et le premier nœud
pris a droite et a gauche de ce point ne peuvent etre egals que dans le seul
cas ou le stilet \struck{est} \struck{ap} supposé appliqué au milieu de la
longueur de la lame dont on cherche a déterminer la maniere de vibrer.
\note{le nombre 130 est à l'encre en haut à droite ; au milieu du haut, un 23 souligné. La première ligne continue la vue 264. L'orthographe de l'autographe est conservée : « égals » et « egals » pour égales, mais « egales » à « a distance egales de part et d'autre ». Le bas du feuillet est blanc sous « vibrer. ».}
\page{268}

14. Pour demontrer cette proposition il faut se rappeller que pour un point
quelconque $L$ situé aux distances $\lambda a$ et $(1-\lambda)a$ des
extremités de la lame, la valeur de l'ordonnée $y$ est en mettant les valeurs
convenables de $\alpha, \beta, \gamma$ et $\delta$, tant pour \struck{le cas
des extremités} fixes que pour le cas celui des extremités libres
\[ C\sin\Big(\zeta+\frac{\omega\omega c\surd 2gb}{aa}\Big)\Big\{e^{\lambda\omega}\pm e^{(1-\lambda)\omega}\ \ill{}\ \{(1\mp e^{\omega})\sin\lambda\omega+(1\pm e^{\omega})\cos\lambda\omega\}\Big\} \]
les signes superieurs étant pris, lorsque les extremités \struck{sont libres}
pour \struck{le cas ou} lorsqu'il y a un nombre impair de nœud, \struck{et
lorsque les extremités sont fixes pour le cas ou il y a un nombre impair} et
les signes inferieurs pour les cas contraires.

cela posé si on veut que le stilet soit appuié au point $L$ on aura pour ce
point $y=0$ c'est adire
\[ e^{\lambda\omega}\pm e^{(1-\lambda)\omega}+(1\mp e^{\omega})\sin\lambda\omega+(1\pm e^{\omega})\cos\lambda\omega=0 \]
et les ordonnées des points situés a droite et a gauche du point $L$ a la
distance $\lambda'a$ de ce point seront
\[ \ldots\ \{e^{(\lambda-\lambda')\omega}\pm e^{(1-\lambda+\lambda')\omega}+(1\mp e^{\omega})\sin(\lambda-\lambda')\omega+(1\pm e^{\omega})\cos(\lambda-\lambda')\}=y \]
\[ \ldots\ \{e^{(\lambda+\lambda')\omega}\pm e^{(1-\lambda-\lambda')\omega}+(1\mp e^{\omega})\sin(\lambda+\lambda')\omega+(1\pm e^{\omega})\cos(\lambda+\lambda')\}=-y \]
les ordonnées sont de signes contraires parceque il est clair que si l'une
est située au dessus de l'axe l'autre doit etre au dessous de la même ligne :
leur somme est donc zéro si elles sont egales entr'elles et on doit avoir
\begin{align*}
0={}& \{e^{\lambda\omega}\pm e^{(1-\lambda)\omega}\}\{e^{\lambda'\omega}+e^{-\lambda'\omega}\}+(1\mp e^{\omega}).2\sin\lambda\omega\,\uncertain{\cos}\lambda'\omega \\
&+(1\pm e^{\omega}).2\cos\lambda\omega\cos\lambda'\omega
\end{align*}
et a cause de $e^{\lambda\omega}\mp e^{(1-\lambda)\omega}=-\{(1\mp e^{\omega})\sin\lambda\omega+(1\pm e^{\omega})\cos\lambda\omega\}$
cette condition devient
\[ \{e^{\lambda\omega}\mp e^{(1-\lambda)\omega}\}\{e^{\lambda'\omega}+e^{-\lambda'\omega}-2\cos\lambda'\omega\}=0 \]
et comme lorsque les extremités sont supposées fixes la seule difference est
entre les valeurs des coëfficiens $\alpha, \beta, \gamma$ et $\delta$
comparés a celles qui sont relatives a l'hypothese des extremités libres
\note{le nombre 131 est à l'encre en haut à droite ; au milieu du haut, un 24 souligné. Le « 14 » qui ouvre la page est dans la marge de gauche, comme le « 13 » de la vue 264. « le cas » est ajouté au-dessus de la ligne devant « celui des extremités libres », et « lors » au-dessus de « le cas ou » biffé. Dans la première formule, le signe qui sépare $e^{(1-\lambda)\omega}$ de l'accolade suivante est noyé sous une rature épaisse et n'est pas lu ; à la formule suivante, la même place porte un « $+$ ». Les doubles signes sont écrits en deux étages, un trait sous ou sur la croix ; ils sont transcrits $\pm$ et $\mp$ comme ils se lisent. Les lignes de points qui ouvrent deux des formules sont de la main de l'auteur et marquent le facteur laissé de côté, comme aux vues 186 et 189 du lot 10. Dans la grande formule, « $\cos$ » devant $\lambda'\omega$ est empâté par une surcharge. Le feuillet s'arrête sur « libres » ; la phrase se poursuit à la vue 270.}
\page{270}

est que les deux dernieres savoir $\gamma$ et $\delta$ sont prises avec un
signe contraire il en resulte seulement que l'on a
\[ e^{\lambda\omega}\pm e^{(1-\lambda)\omega}-\{(1\mp e^{\omega})\sin\lambda\omega+(1\pm e^{\omega})\cos\lambda\omega\}=0 \]
\[ \ldots\ \{e^{(\lambda-\lambda')\omega}\pm e^{(1-\lambda+\lambda')\omega}-\{(1\mp e^{\omega})\sin(\lambda-\lambda')\omega+(1\pm e^{\omega})\cos(\lambda-\lambda')\omega\}=y \]
\[ \{e^{(\lambda+\lambda')\omega}\pm e^{(1-\lambda-\lambda')\omega}-\{(1\mp e^{\omega})\sin(\lambda+\lambda')\omega+(1\pm e^{\omega})\cos(\lambda+\lambda')\omega\}=-y \]
au lieu des valeurs précédantes ce qui n'empêche pas de conclure
\[ \{e^{\lambda\omega}\mp e^{(1-\lambda)\omega}\}\{e^{\lambda'\omega}+e^{-\lambda'\omega}-2\cos\lambda'\omega\}=0 \]
On voit aisément qu'en supposant $\lambda=\frac{1}{2}$ cette condition est
toujours remplie car il faut alors prendre le signe superieur dans le
facteur $e^{\lambda\omega}\mp e^{(1-\lambda)\omega}$ puisque dans ce cas il y a
necessairement un nombre impair de nœuds et il est clair que
$e^{\frac{1}{2}\omega}-e^{(1-\frac{1}{2})\omega}=0$ ; mais on voit en même tems
que cette valeur de $\lambda$ est la seule pour laquelle ce facteur puisse
etre egalé a zéro et que dans tout autre cas on mené a conclure
$e^{\lambda'\omega}+e^{-\lambda'\omega}=2\cos\lambda'\omega$ resultat absurde pour tout
une valeur quelconque de $\lambda'$ differente de zéro.

Ainsi il est prouvé a plus forte raison qu'un nœud ne peut jamais etre
placé a distance egals des nœuds les plus voisins pris a droite et a gauche
car il faudroit que les ordonnées \struck{fussent} de ces derniers nœuds
fussent égals a zéro, et par conséquent aussi leur somme, ce que l'on vient
de demontrer être impossible pour un nœud situé dans un lieu quelconque
different du milieu d'une lame dont les extremités sont supposées libres ou
fixes $+$

15. $*$ Au contraire lorsqu'il s'agit d'une lame dont les deux extremités
sont simplement appuiées si on veut que le stilet soit appuié \struck{de
maniere} en un point $L$ pris a distance $\lambda a$ d'une des extremités et
qu'il existe deux autres nœuds aux distances $(\lambda-\lambda')a$ et
$(\lambda+\lambda')a$ de la même extremité de sorte que \struck{que le nœud}
\note{le nombre 132 est à l'encre en haut à droite ; au milieu du haut, un nombre souligné dont les deux chiffres sont repassés l'un sur l'autre, lu 25. Le « 15 » qui ouvre le dernier alinéa est dans la marge de gauche. Dans la deuxième formule, le signe qui précède l'accolade et celui qui précède $y$ sont deux taches épaisses recouvrant ce qui est écrit ; ils sont rendus « $-$ » et « $=$ » d'après la troisième formule, qui est la même à un signe près, et le sinus porte au milieu une surcharge. Le « $+$ » qui ferme l'alinéa 14 et le « $*$ » qui ouvre l'alinéa 15 sont les signes de renvoi des deux notes marginales de la vue 264. Le feuillet s'arrête sur un mot biffé ; la phrase se poursuit à la vue 272.}
\page{272}

$*$ Je ne crois pas que l'on ait encore fait cette remarque de l'inégalité
\struck{entre} des parties comprises entre \struck{deux} nœuds de vibration ;
car Mr Chladni se contente de dire \struck{qu'un} en partant de la lame dont
les deux extremités sont libres qu'une partie comprise entre deux nœuds est a
peu près double d'une partie située a une extremité ce qui semble
\struck{supposer} que est a peu près moitié moins longue qu'une partie
comprise entre deux nœuds, ce qui semble supposer que toutes les parties
soumises a cette derniere condition \struck{sont} de la même longueur.
\uncertain{V.} traité d'acoustique pag 4 §§ 71

Cependant Mr Ricati dans le memoire que j'ai déja cité, ouvrage dans lequel
il a entrepris developper l'analyse du cas des deux extremités libres, en
suivant la methode d'Euler, et où il a calculé avec plus de soins les valeurs
qui entrent dans \struck{les} formules donne pour \struck{la quantité} que
j'ai jusqu'a été designée ici par $\omega$ des valeurs qui comparés
entr'elles menent a conclure que pour le cas des deux extremités libres une
partie située a l'extremité est moins grande que la moitié de la partie
comprise entre les deux nœuds \struck{placés} les plus voisins de cette
extremité, que la partie comprise entre le second nœud et le troisieme est
plus grande que celle qui est située entre le premier et le second, que de
même encore la partie qui \struck{se trouveroit} comprise entre le troisieme
et le quatrieme est plus grande que toutes les précédantes \struck{et toute}
qu'en général une partie comprise entre deux nœuds est d'autant plus longue
qu'elle est plus éloignée de l'extremité et plus voisine du milieu de la
lame. et que cependant les differences sont beaucoup plus grandes entre les
parties qui sont pres des extremités qu'entre celles qui s'en eloignent
d'avantage

2. Mais Mr Ricati se contente de donner les valeurs dont j'ai parlé sans en
conclure l'inegalité des parties dont il s'agit. V. §§ XXX p.
\uncertain{503}. $*$
\note{le nombre 133 est à l'encre en haut à droite ; au milieu du haut, un nombre souligné dont le second chiffre est repassé, lu 25, précédé à gauche, en petits caractères, d'une mention lue « 2.\textsuperscript{eme} » avec doute. La lettre capitale suivie d'un point qui précède « traité d'acoustique » est tracée d'une seule boucle : elle est lue « V. » (voyez) d'après le « V. » de la dernière ligne, mais la forme n'est pas la même. « et que » est ajouté au-dessus de la ligne devant « cependant », « les » au-dessus de la ligne devant « parties soumises », « deux » au-dessus de la ligne dans « des deux extremités libres ». Les deux dernières lignes sont d'une écriture beaucoup plus fine, ajoutées au bas du feuillet ; dans leur renvoi, trois croix épaisses tiennent la place du numéro des paragraphes, et le nombre de la page, lu 503, a un chiffre du milieu qui peut être un zéro ou un o. Le « $*$ » qui ouvre la page et celui qui ferme la dernière ligne sont des signes de renvoi.}
\page{274}

point d'appui du stilet equivalent à un nœud soit placé a distance egale de
deux autres nœuds situés a droite et a gauche ces differentes conditions
seront exprimées par les équations
\[ \sin\lambda\omega=0, \qquad \sin(\lambda-\lambda')=0 \]
\[ \sin(\lambda+\lambda')=0 \]
et on \struck{conclut} trouve en ajoutant les deux dernieres
$2\sin\lambda\cos\lambda'\omega=0$ qui est evidenment satisfaite \struck{pour
toute valeur convenable a $\omega$ et independam} independamment
de la valeur de $\lambda'$ en vertu de la premiere de ces équations.

On voit donc a present par quelle raison le cas des deux extremités appuiées
est le seul pour lequel on puisse avoir une solution générale \struck{de} des
équations trouvées en suivant la methode d'Euler \struck{\ill{}}

\marginal{$\times$ Dans les \uncertain{problemes} ou il s'agit de déterminer
l'effet de l'application d'un stilet $\times$}

16. On \struck{voit} peut remarquer. \struck{qu'} $\times$ il existe la plus
grande analogie entre les équations qui \struck{determinent} servant a
trouver \struck{\ill{}} la valeur de l'angle $\omega$ pour \struck{dans} les
differens cas relatifs a l'état des extremités et aussi entre les facteurs de
ces équations et les conditions par lesquelles l'angle $\omega$ est determiné
\struck{dans} lorsque l'on a simplement egard a l'état des extremités, sans
supposer l'application du stilet dans aucun \struck{point} de l'étendue des
lames dont on veut déterminer le mouvement

Afin que l'on puisse saisir plus aisement les rapports entre ces differentes
\struck{quantités} equations je vais les \struck{ray} écrire ici en y
joignant celle que j'ai également calculés pour le cas ou le stilet
\note{le nombre 134 est à l'encre en haut à droite ; au milieu du haut, un 26 souligné. Le « 16 » de l'alinéa est dans la marge de gauche, en gros chiffres empâtés. La note marginale est de la main de l'auteur et appelée par une croix placée dans le texte après « qu' ». « trouve en ajoutant » est écrit au-dessus de « conclut » biffé ; « peut remarquer. » au-dessus de « voit » biffé ; « servant a trouver » et les mots biffés qui le suivent, au-dessus de « determinent » biffé, et ces mots biffés ne sont pas lus ; « pour » au-dessus de « dans » biffé ; « je vais les ray écrire ici » au-dessus de la ligne, entre « equations » et « en y joignant », et mis à sa place, « ray » y étant biffé. Le feuillet s'arrête sur « le stilet » ; la phrase se poursuit à la vue 276.}
\page{276}

et supposé appliqué dans un point quelconque de l'étendue d'une lame
elastique vibrante dont une extrémité seroit fixée et l'autre simplement
appuiée.

Voici donc d'abord les équations qui déterminent la valeur de l'angle
$\omega$ dans les differens cas ou on a seulement égard a l'état des
extremités :

\begin{itemize}
\item Pour les deux extremités libres. cas 1.\textsuperscript{er} N\textsuperscript{o} 9 (memoire d'Euler) $\ldots$ $\cos\omega=\frac{2e^{\omega}}{e^{2\omega}+1}$
\item Pour une extremité libre et l'autre appuiée. cas 2.\textsuperscript{eme} N\textsuperscript{o} 23 $\ldots$ $\text{tang}\,\omega=\frac{e^{2\omega}-1}{e^{2\omega}+1}$
\item Pour une extremité libre et l'autre fixée. cas 3.\textsuperscript{eme} N\textsuperscript{o} 30 $\ldots$ $\cos\omega=-\frac{2e^{\omega}}{e^{2\omega}+1}$
\item Pour les deux extremités appuiées. cas 4.\textsuperscript{eme} N\textsuperscript{o} 35 $\ldots$ $\sin\omega=0$
\item Pour une extremité appuiée et l'autre fixée cas 5.\textsuperscript{eme} N\textsuperscript{o} 39 $\ldots$ $\text{tang}\,\omega=\frac{e^{2\omega}-1}{e^{2\omega}+1}$
\item Pour les deux extremités fixées cas 6.\textsuperscript{eme} N\textsuperscript{o} 43 $\ldots$ $\cos\omega=\frac{2e^{\omega}}{e^{2\omega}+1}$
\end{itemize}

On \struck{peut} voit déja \struck{remarquer} qu'il existe les plus grands
\struck{analogies} rapports, entre l'effet de la fixation et de l'entiere
liberté d'une extremité car on voit que l'on peut changer ces conditions
l'une en l'autre sans que l'équation qui détermine la valeur de l'angle
$\omega$ change de forme.

Apresent, les équations qui déterminent la valeur de l'angle $\omega$,
lorsque l'on a egard a l'application du stilet a un point quelconque de
l'étendue de la lame elastique vibrante sont pour le cas des deux extremités
appuiées :
\[ \sin\lambda\omega\Big\{\frac{e^{(1-\lambda)\omega}+e^{(\lambda-1)\omega}}{e^{(1-\lambda)\omega}-e^{(\lambda-1)\omega}}\sin(1-\lambda)\omega-\cos(1-\lambda)\omega\Big\}+\sin(1-\lambda)\omega\Big\{\frac{e^{\lambda\omega}+e^{-\lambda\omega}}{e^{\lambda\omega}-e^{-\lambda\omega}}\sin\lambda\omega-\cos\lambda\omega\Big\}=0 \]
Pour les deux extremités libres
\begin{align*}
&\Big\{\frac{2}{e^{\lambda\omega}+e^{-\lambda\omega}}+\cos\lambda\omega\Big\}\Big\{\frac{e^{(1-\lambda)\omega}+e^{(\lambda-1)\omega}}{e^{(1-\lambda)\omega}-e^{(\lambda-1)\omega}}\sin(1-\lambda)\omega-\cos(1-\lambda)\omega\Big\} \\
&+\Big\{\frac{2}{e^{(1-\lambda)\omega}+e^{(\lambda-1)\omega}}+\cos(1-\lambda)\omega\Big\}\Big\{\frac{e^{\lambda\omega}+e^{-\lambda\omega}}{e^{\lambda\omega}-e^{-\lambda\omega}}\sin\lambda\omega-\cos\lambda\omega\Big\}=0
\end{align*}
Pour les deux extremités fixées
\begin{align*}
&\Big\{\frac{2}{e^{\lambda\omega}+e^{-\lambda\omega}}-\cos\lambda\omega\Big\}\Big\{\frac{e^{(1-\lambda)\omega}+e^{(\lambda-1)\omega}}{e^{(1-\lambda)\omega}-e^{(\lambda-1)\omega}}\sin(1-\lambda)\omega-\cos(1-\lambda)\omega\Big\} \\
&+\Big\{\frac{2}{e^{(1-\lambda)\omega}+e^{(\lambda-1)\omega}}-\cos(1-\lambda)\omega\Big\}\Big\{\frac{e^{\lambda\omega}+e^{-\lambda\omega}}{e^{\lambda\omega}-e^{-\lambda\omega}}\sin\lambda\omega-\cos\lambda\omega\Big\}=0.
\end{align*}
\note{le nombre 135 est à l'encre en haut à droite ; au milieu du haut, un 27 souligné. « d'abord » est ajouté au-dessus de la ligne après « Voici donc », « appuiées » au-dessus de la ligne à la fin de « pour le cas des deux extremités », « voit » au-dessus de « peut » biffé et « rapports » au-dessus de « analogies » biffé. Les six renvois sont à la suite des lignes, séparés du membre de droite par une ligne de points de conduite ; ils sont donnés ici en liste. Les numéros d'ordre des cas sont écrits d'une main où le 3 se ferme comme un 9 et le 6 s'ouvre comme un b ; ils se lisent 1, 2, 3, 4, 5, 6 et les numéros de paragraphe 9, 23, 30, 35, 39, 43. La dernière formule de la page est celle qui ouvre la vue 262.}
\page{278}

En employant les équations resultantes de l'état des extremités et les
combinant avec les quatre equations données par la condition de
l'application du stilet \marginal{voyez § 47 et suiv. du memoire d'Euler}
j'ai encore trouvé pour une extremité fixée et l'autre libre
\begin{align*}
&\Big\{\text{tang}\,\lambda\omega-\frac{e^{\lambda\omega}-e^{-\lambda\omega}}{e^{\lambda\omega}+e^{-\lambda\omega}}\Big\}\Big\{\text{tang}(1-\lambda)\omega-\frac{e^{(1-\lambda)\omega}-e^{(\lambda-1)\omega}}{e^{(1-\lambda)\omega}+e^{(\lambda-1)\omega}}\Big\} \\
&=2\,\text{tang}\,\lambda\omega\frac{e^{\lambda\omega}-e^{-\lambda\omega}}{e^{\lambda\omega}+e^{-\lambda\omega}}\Big\{1-\frac{2}{(e^{(1-\lambda)\omega}+e^{(\lambda-1)\omega})\cos(1-\lambda)\omega}\Big\}
\end{align*}

17. Il est aisé de remarquer que dans le cas ou \struck{\ill{}} l'état des
extremités \struck{est} semblable, ces équations sont composées de deux
termes dont l'un est le produit de la quantité qui devroit etre egalé a
zero, si la portion de la lame de longueur $\lambda a$, se mouvoit comme une
lame entière et separée de même longueur, et dont une extremité seroit
appuiée au point d'application du stilet ; tandis que l'autre \struck{seroit}
resteroit dans l'état déterminé pour les extremités de la lame entière ; par
la quantité qui devroit etre egalé a zéro si la portion de longueur
$a(1-\lambda)$, se mouvoit comme une lame pareillement separée et dont une
extremité seroit fixée au point d'application du stilet tandis que l'autre
resteroit dans l'état déterminé pour les extremités de la lame entière
\struck{et} que l'autre terme est pareillement composé des mêmes quantités en
changeant $\lambda a$ en $(1-\lambda)a$ et vice versa.

On peut encore observer que même dans le cas ou l'état des extremités n'est
plus semblable il existe \struck{encore} une grande analogie entre l'équation
qui y est relative et celle des cas precedans car si on considere le stilet
comme simplement appuié on a \struck{a la fois $\text{tang}\,\lambda\omega=0$}
ou l'une des équations
\[ \sin\lambda\omega=0 \quad\text{et}\quad \text{tang}(1-\lambda)\omega=\frac{e^{(1-\lambda)\omega}-e^{(\lambda-1)\omega}}{e^{(1-\lambda)\omega}+e^{(\lambda-1)\omega}} \]
\note{le nombre 136 est à l'encre en haut à droite ; au milieu du haut, un 28 souligné, dont le second chiffre est empâté. Le « 17 » de l'alinéa est dans la marge de gauche. Le renvoi « voyez § 47 et suiv. du memoire d'Euler » est écrit en petits caractères à la suite de la ligne, dans la marge de droite ; il désigne le Problema de la vue 193 du lot 10. Le mot biffé après « dans le cas ou » est recouvert d'une rature épaisse et n'est pas lu. « équations » est ajouté au-dessus de la ligne à la fin de « ou l'une des ». La ligne biffée « a la fois tang $\lambda\omega=0$ » est barrée d'un long trait horizontal.}
\page{280}

\struck{l'une de ces équations} est la condition du mouvement d'une lame de
longueur $\lambda a$ dont \struck{les} une des extremités \struck{seroit}
appuiée au point d'application du stilet et dont l'autre seroit également
appuiée en vertu de la \struck{supp} l'état des extremités de la lame
primitive, l'autre est la condition du mouvement d'une lame de longueur
$(1-\lambda)a$, dont une des extremités seroit appuiée au point d'application
du stilet et dont l'autre seroit fixée en vertu de l'état des extremités de
la lame primitive

et de même si on considere le stilet comme fixé l'une des équations
\[ \text{tang}\,\lambda=\frac{e^{\lambda\omega}-e^{-\lambda\omega}}{e^{\lambda\omega}+e^{-\lambda\omega}},\qquad
\cos(1-\lambda)\omega=\frac{2}{e^{(1-\lambda)\omega}+e^{(\lambda-1)\omega}} \]
et la condition d'un mouvement d'une lame de longueur $\lambda a$ dont une
des extremités seroit fixée au point d'application du stilet et dont l'autre
seroit appuiée en vertu de l'état des extremités de la lame primitive tandis
que l'autre est la condition du mouvement d'une lame de longueur
$(1-\lambda)a$ dont les deux extremités seroient fixées.

\marginal{la methode qui conduit a tous ces resultats est loin d'être aussi
utile qu'on auroit pû d'abord l'esperer, car elle semble indiquer des
conclusions qui ne sont pourtant pas applicables et les seules
\struck{consé}quences admissibles que l'on puisse en déduire \struck{des
équations} \struck{tout} \uncertain{dedu} sont d'abord}

18. Cependant malgré des analogies si remarquables cette méthode ne mene a
aucun resultat applicable puisque, comme je l'ai montré ci dessus §.
\uncertain{9} excepté pour le cas ou les deux extremités sont appuiées il
n'y a pas de solution générale de ces equations et que par conséquent on ne
peut, jamais pour aucun pour le cas ou les extremités sont appuiées de
montrer q\ill{} si il y a un point d'appui aux distances $\lambda a$ et
$(1-\lambda)a$ des extremités de la lame on devra avoir $\sin\lambda\omega=0$,
et qu'alors il se formera des points de repos a des distances $2\lambda a$ et
$(1-2\lambda)a$, $3\lambda a$ et $(1-3\lambda)a$ \&c. des mêmes extremités
parceque la condition \struck{est $\sin\lambda\omega=$} entraine les
suivantes $\sin 2\lambda\omega=0$, $\sin 3\lambda\omega=0$ \&c.
\note{le nombre 137 est à l'encre en haut à droite ; au milieu du haut, un 29 souligné, dont le second chiffre est repassé. Le « 18 » de l'alinéa est dans la marge de gauche, en gros chiffres empâtés. « dont » est ajouté au-dessus de la ligne après « du stilet et ». Dans la première des deux équations de la seconde partie, la tangente est écrite « tang. $\lambda$ » sans $\omega$, alors que les autres termes portent $\lambda\omega$ ; la lecture est laissée telle. La note marginale de gauche reprend l'alinéa 18 ; ses deux dernières lignes sont surchargées et le mot qui précède « sont d'abord » n'est pas assuré. Un long trait de plume oblique, épaissi de taches, traverse l'alinéa 18 depuis « cette méthode » jusqu'à « on ne peut » ; il n'est pas dit ici s'il annule le passage. Après « de montrer q », le feuillet porte un petit manque de papier et la fin du mot n'est pas lue. La page se termine sur « \&c. » ; la suite est à la vue 281, hors de ce lot.}

\end{document}
